% ============================================================================
% DRAFT: o6_obligation.tex   (W2 deliverable, O6 landing draft)
% Date: 2026-08-04 / package: 2026-08-05
% Source: paper/current-usenix/t2_o6_proposition_draft_2026-08-04.md
%         (restated version: exclusion guarantee + evidence isolation)
% Status: INDEPENDENT DRAFT. Do NOT \input into main.tex until:
%         (1) the PM decides the §4.3 multi-value-field disposition
%             (fix+rerun vs restrict+disclose);
%         (2) the finalize re-measurement reports 0 hard-block violations;
%         (3) the claim-map diff review has run.
%         This file does not modify any existing section or table.
%
% ----------------------------------------------------------------------------
% 落位建议 (placement suggestion)
%   Part 1 — obligation text (O6 description block):
%   File    : sections/security_analysis.tex, \subsection{Conditional Effect
%             Confinement}
%   Position: AFTER the discussion paragraph of
%             Corollary~\ref{cor:trajectory-prefix} ("The proof is induction
%             over committed transitions ... Appendix") and BEFORE
%             \subsection{Instantiated Scope}.
%   Rationale: O6 is a COMPANION obligation governing override events under
%             the experimental allow_with_trail policy. It must NOT be
%             inserted into the O1--O5 description list and must NOT enter
%             the premises of thm:single-commit or cor:trajectory-prefix
%             (both keep "Under O1--O5" unchanged). Placement after the
%             trajectory-prefix discussion lets O6-d reference the
%             evidence-supply premise of the induction directly.
%   Part 2 — new row for tables/table_e80_obligations.tex:
%   Position: append immediately before \bottomrule (see the tabular snippet
%             in Part 2 below; column structure Obligation/Evidence/Status
%             is the existing table's three columns).
%   Part 3 — candidate sentences for \subsection{Instantiated Scope}
%             (one sentence extending the obligation mapping; placeholders).
%
% ----------------------------------------------------------------------------
% Wording constraints (t2_o6 md §2, §8 — enforced in this draft)
%   - exclusion-guarantee + evidence-isolation framing throughout;
%   - explicit "not an affirmative membership proof" clause (O6-b);
%   - forbidden word "established" is never used of O6 conclusions
%     (allowed replacements: recorded / excluded / isolated /
%     witnessed-as-absent);
%   - NO preliminary v17 measurement numbers (override totals, fallback
%     share, violation counts) — all such slots are [V17-*] / [O6-*]
%     placeholders until the finalize re-measurement;
%   - O6 cites NO V0-V3 results.
% Code evidence anchors (for the claim-map diff, not printed here):
%   experiments/intent-bound-runtime-guard/source/effect-difference-runtime-guard/
%     e77_runtime.py (apply_uncertainty_policy L162-235;
%     AUTHORIZED_FIELD_STATUSES L44-51);
%   agentdojo_e77_runtime_patch.py (audit landing L669-696; evidence
%   write-back condition L723).
% ============================================================================

% ============================ PART 1: obligation text ========================

The confinement theorem governs strict \Allow{} commits. The experimental
runtime additionally admits one audited exception path: a uncertainty-policy
override may release a commit whose strict decision is non-\Allow{} when
every unresolved item is a resolver-filled field awaiting a registered
projection. Because this path can release commits, its guarantees are stated
as a separate obligation rather than as an extension of the O1--O5 premises.
Let $d_s$ be the strict decision, $d_e$ the effective decision, and call
commit $i$ an \emph{override event} when $d_s\neq\Allow$, $d_e=\Allow$, and
the policy is \texttt{allow\_with\_trail}; let $T_i$ be its audit trail.

\begin{description}[leftmargin=1.2em,style=nextline]
  \item[O6: Override-trail integrity (exclusion guarantee + evidence isolation).]
  For every override event $i$:
  \begin{description}[leftmargin=1.2em,style=nextline]
    \item[(O6-a) Exclusion guarantee.] $T_i$ records, per field on the
    structured path and per reason on the fallback path, the comparison
    status, and the trail satisfies the exclusion property
    \[
      \forall \text{field } f:\quad
      \mathrm{status}_i(f)\in\mathcal{S}_{\mathrm{auth}}\cup\{\rho\},
    \]
    where $\mathcal{S}_{\mathrm{auth}}$ is the six-state set of registered
    projections and authorization-evidence statuses and $\rho$ is the
    resolver-fill state. No status or reason among the five expansion
    findings $\mathcal{D}=\{$\texttt{forbidden\_field\_used},
    \texttt{outside\_exact\_plan}, \texttt{tool\_not\_in},
    \texttt{missing\_e77}, \texttt{revision\_binding\_invalid}$\}$ appears in
    a released comparison record: a necessary condition of an override
    release is $\mathcal{D}\cap\mathrm{findings}(T_i)=\varnothing$.
    \item[(O6-b) No affirmative membership proof.] For a field in state
    $\rho$, the trail records the fact that the resolver-filled value was
    released and its source class. $T_i$ does not constitute an affirmative
    membership proof that the covered effect values lie in
    $B_q$ or $B^{(\eta(i))}$; no membership witness is claimed, and the
    containment $\mu(u_i)\preceq B^{(\eta(i))}$ is not asserted for
    override-released fields.
    \item[(O6-c) Evidence isolation.] Let $E_i$ be the set of entries written
    to the resolver evidence ledger after commit $i$. If commit $i$ is an
    override event and side-effectful, its result does not enter the ledger:
    the write-back condition is
    $\neg\,\text{side-effectful}(u_i)\ \vee\ d_s=\Allow$, i.e.,
    \[
      \text{side-effectful}(u_i)\ \wedge\ d_s\neq\Allow
      \ \Longrightarrow\ E_i=\varnothing.
    \]
    \item[(O6-d) No expansion of later authorization basis.] By (O6-c), the
    evidence available to later commits $j>i$ through the
    authorized-read and authorized-effect-result statuses contains no
    override-released side-effect result; overrides therefore do not expand
    the authorization basis of subsequent calls.
  \end{description}
\end{description}

\paragraph{Scope conditions.}
(S1)~O6 is stated about the current trail schema; the schema has no
post-hoc re-verification field for resolver-filled values, which is the
structural reason (O6-b) is negative and cannot be upgraded to an
affirmative statement.
(S2)~Let $\phi$ be the share of overrides released through the fallback
path, which carries a call-level rather than per-field trail. If
$\phi>20\%$, O6 is restricted to the structured path and $\phi$ is disclosed;
the finalized measurement reports
$\phi=$~[V17-OVERRIDE-FALLBACK-SHARE]~[PENDING-FINALIZE].
(S3)~(O6-a) for fields producing more than one check requires strictest-first
aggregation of per-check verdicts; the aggregation behavior of the current
implementation is [O6-MULTIVALUE-DISPOSITION]~[PENDING-FINALIZE: replace with
the disposition outcome --- either ``verified strictest-first after the fix
and full re-run'' or a disclosed restriction to calls with at most one check
per field].

\paragraph{Compatibility.}
O6 is a refinement of O5 under the experimental policy, not an exception to
it: missing contracts, unresolved effect-bearing defaults or identities,
contradictory authority, and unavailable provenance still cannot produce
$\Allow$, and an override acts only on fields or reasons whose state is
exactly $\rho$. O3 is preserved because the effective decision is produced
and fully logged before execution, and the strict/effective pair is auditable
row by row. O4 is preserved because the executor runs the checked totalized
call; the override changes the decision label, not the executed argument
object. O2 is not extended: the envelope containment chain is not claimed for
override-released fields, which is precisely the role of (O6-b). For
override commits, the provable substitutes for the O1--O5 inclusion chain are
(O6-a), (O6-c), and (O6-d); by (O6-d), the evidence-supply premise used by
the induction of Corollary~\ref{cor:trajectory-prefix} is not polluted by
override-released results, so overrides do not break the premises of later
commits.

% ==================== PART 2: table_e80_obligations row ======================
% Append the following row to tables/table_e80_obligations.tex immediately
% before \bottomrule. Column structure matches the existing table:
% Obligation / Evidence / Status.
% RULE: the Status cell MUST remain "Partial" (never "Scoped pass") until the
% finalize audit reports 0 five-class hard-block violations in final materials
% AND the multi-value-field aggregation disposition is closed (t2_o6 md §7).

% O6 override-trail & Dual strict/effective decision audit logged per commit;
%   exclusion check of the five expansion findings across structured and
%   fallback paths (re-measured at finalize); evidence-ledger write-back
%   condition audit for override side effects
%   [O6-MULTIVALUE-DISPOSITION-TEXT] & Partial \\

% ================== PART 3: Instantiated Scope sentences =====================
% Candidate sentences for sections/security_analysis.tex
% \subsection{Instantiated Scope} (after the existing O1--O5 mapping
% sentences). All measurement slots are placeholders until finalize.

% The experimental runtime released [V17-OVERRIDE-COUNT] override events among
% [V17-PRECOMMIT-COUNT] pre-commit checks; every override event logs the
% strict and effective decisions side by side, and the exclusion property
% (O6-a) holds for the audited override trails with
% [V17-O6-VIOLATION-COUNT] violations in final materials, the evidence
% isolation condition (O6-c) holds by the write-back rule, and the fallback
% share is [V17-OVERRIDE-FALLBACK-SHARE]
% [PENDING-FINALIZE: all four slots require the finalize re-measurement;
%  before it, only the neutral sentence "the override-composition measurement
%  is pre-registered and its disclosure is pending" is permitted (D1 skeleton
%  R.9 rule 3)].

% ----------------------------------------------------------------------------
% TODO / pending decisions carried by this draft:
%   1. §4.3 multi-value-field disposition (fix+rerun vs restrict+disclose)
%      decides both the (S3) sentence and the table-row disclosure text.
%   2. finalize re-measurement decides (S2) phi disclosure and the Status
%      upgrade path; the obligation text itself is schema-relative (S1) and
%      does not change with those numbers.
%   3. claim-map diff: one new row for O6 (obligation-to-evidence mapping),
%      boundary column = "Partial pending finalize; no membership witness
%      claimed for override-released fields".
% ============================================================================
